\section{Disintegration into original ambient currents}
\label{sec:r12-original-current-disintegration}

The closed parameter bands admit a genuine local disintegration,
but its conditional currents remain on the original ambient
$K$-analytic space. They are not currents on a newly introduced
residue-field analytic space. This distinction is essential for
using the frozen arbitrary-background BT theorem without changing
its field or its test sheaf.

Throughout this section $U$ is second countable in addition to the
standing original hypotheses. The conclusion is ordinary-local on
every space in the project: every point has an ordinary neighborhood
of this kind, as proved below. Let $s:U\to J\subset\mathbf R^c$ be
the original parameter map, and let $R$ be an allowed $s$-background.

\subsection{A countably cofinal system of original smooth tests}

\begin{lemma}[Local countability does not discard actual refinements]
\label{lem:r12-countably-cofinal-original-tests}
Every point of an allowed original analytic space has a second
countable ordinary neighborhood. On any such $U$, the original
compact smooth CLD test spaces in all bidegrees admit a countable
cofinal family of finite moment-and-cutoff presentation stages.
Each stage is a quotient of a separable smooth coefficient space,
has a countable dense subset, and has a countable dense subset in
its original middle-positive cone. All equality kernels between
the stages can likewise be checked on countably many dense
coefficient relations.

More precisely, every original compact test, including a test
initially represented after an arbitrary actual refinement, is a
finite sum of tests in these stages. Approximation within each
stage is in its usual smooth coefficient topology on one fixed
compact support. No assertion about all GK delta-tests is made.
\end{lemma}

\noindent\textbf{Status: \texttt{NEW\_PROVED}} for the exact
original-test cofinality used in the conditional-current argument.
The original smooth cutoff/partition and common-moment properties
are frozen \texttt{INTERNAL\_INPUT}s.

\begin{proof}
The finite extension $K/\mathbf Q_p$ is separable as a metric field.
Every $K$-affinoid algebra is a quotient of a generalized Tate
algebra in finitely many variables. Polynomials with coefficients
in a fixed countable dense subfield, followed by passage to the
quotient, give a countable dense subset of that Banach algebra.
This argument also applies to non-strict radii: finite polynomials
are dense in the defining weighted completion.

The Berkovich spectrum of a separable affinoid algebra is compact
metrizable. Indeed a countable dense set of functions determines
each bounded multiplicative seminorm, and all remaining evaluations
are uniform limits of these evaluations, since
$\bigl||f(x)|-|g(x)|\bigr|\le\|f-g\|$. The spectrum embeds
continuously and injectively in a countable product of compact
real intervals. Compactness and the Hausdorff property identify
its topology with that induced metric topology. By goodness, an
ordinary point has an affinoid neighborhood; its ordinary interior
is the required second countable neighborhood.

On a second countable good $U$, choose a countable collection of
affinoid neighborhoods whose ordinary interiors form a basis.
To construct it, for each member of a countable topological basis
cover its points by affinoid interiors contained in that member,
and choose a countable subcover. Choose a countable dense subset
$D_B$ of each resulting affinoid algebra $\mathcal O(B)$.
Whenever $g$ is an analytic unit on a neighborhood of $B$, choose
$g_0\in D_B$ with
\[
 \|g-g_0\|_B<\min\left\{\inf_{x\in B}|g(x)|,\|g^{-1}\|_B^{-1}\right\}.
 \tag{DI1}
\]
The right side is positive by compactness and invertibility. The
ultrametric inequality gives $|g_0|=|g|$ on $B$, and the Neumann
series applied to $1+g^{-1}(g_0-g)$ makes $g_0$ a unit there. Their actual
logarithmic moment functions and their original differentials
therefore agree exactly on the ordinary interior. Thus every finite
actual moment tuple on a compact local stage is replaceable, after
finite ordinary localization, by a tuple from a countable family.
This applies equally to arbitrary extra units in a refinement.

Fix countably many original smooth cutoffs for pairs of basis
neighborhoods with nested compact closures, equal to one on the
smaller compact and supported in the larger. The frozen cutoff
theorem supplies each such choice. Include the finitely many
moments used by each cutoff in the countable family; finite unions
and finite common refinements remain countable. A finite collection
of these cutoffs covers the support of any test. Normalizing their
sum on a neighborhood of that support yields a finite smooth
partition and reduces the test to compact stages in the stated
countable family.

For a fixed tuple and cutoff, write its coefficient arrays on
real coordinate neighborhoods. Use a countable exhaustion by
rational real balls and compact coordinate cores. The corresponding
$C^\infty$ coefficient spaces are separable metrizable spaces;
finite products and their quotients have the same property. The
actual zero-form conditions on a fixed moment image are closed
linear conditions, obtained by restrictions to its relevant tangent
spaces. The smooth coefficient relations comparing two stages can
be imposed on their countable common refinements; their kernels
are separable subspaces. Choose a countable dense subset in each
such kernel. The preimage of the middle-positive cone is likewise
a subspace subset of a second countable metric coefficient space,
and hence is separable; choose dense positive representatives.
These choices prove all countability assertions. No rationality of
the real coefficients of a test is required: only density within
its fixed smooth presentation is used.
\end{proof}

\begin{lemma}[A recalled Radon disintegration construction]
\label{lem:r12-Radon-conditional-measures}
Let $M$ be a locally compact second countable Hausdorff space,
$s:M\to J$ continuous, and $\mu$ a finite positive Radon measure.
Put $\nu=s_*\mu$. There is a measurable family of positive Radon
probability measures $\mu_a$, for $\nu$-almost every $a$, carried
by $s^{-1}(a)$, such that for every nonnegative Borel $F$,
\[
                       \int_M F\,d\mu
                          =\int_J\left(\int_M F\,d\mu_a\right)d\nu(a).
 \tag{DI2}
\]
The identity also holds after insertion of any bounded Borel
parameter weight. Integrable signed functions are allowed.
\end{lemma}

\noindent\textbf{Status: \texttt{EXTERNAL\_KNOWN}} as the usual
Radon conditional-measure theorem. The construction below records
the countable version needed here, using Radon--Nikodym and Riesz
representation only.

\begin{proof}
Take a countable dense rational vector lattice in $C_0(M)$, enlarged
to include a compact approximate unit $\chi_m\uparrow1$, and products
with pullbacks of a countable determining algebra in $C_c(J)$.
For each chosen $f$, the signed measure
$E\mapsto\int_{s^{-1}(E)}f\,d\mu$ is absolutely continuous with
respect to $\nu$. Choose its Radon--Nikodym derivative $g_f$.
Outside one null set the maps $f\mapsto g_f(a)$ are rational-linear,
positive, bounded by $\|f\|_\infty$, and satisfy
\[
                  g_{\phi(s)f}(a)=\phi(a)g_f(a).
\]
There are only countably many identities and inequalities to impose.
Extend by continuity to $C_0(M)$ and apply Riesz representation to
obtain $\mu_a$ of mass at most one. The approximate unit and
monotone convergence show that its mass equals one almost everywhere:
the increasing limit of $g_{\chi_m}$ has integral $\nu(J)$ and is
bounded above by one. The displayed multiplication identity then
gives $s_*\mu_a=\delta_a$, first on the determining continuous
algebra and then on Borel sets. Thus $\mu_a$ is carried by
$s^{-1}(a)$.

The derivatives $g_f$ supply measurability and (DI2) on the chosen
algebra. Uniform approximation extends these assertions to
$C_0(M)$; the monotone-class theorem extends them to nonnegative
Borel functions. Multiplication by a bounded Borel function of
$s$ gives the weighted version. The signed assertion follows from
positive and negative parts. No regularity of a conditional density
in $a$ is asserted.
\end{proof}

\subsection{Conditional currents retain the entire original test sheaf}

\begin{theorem}[Disintegration of a parameter-annihilated original current]
\label{thm:r12-original-current-disintegration}
Let $R$ be an allowed $s$-background on second countable $U$.
There is a finite positive Radon control measure $\nu$ on $J$ and
a measurable family of original currents $R_a$ on the same $U$
such that, for $\nu$-almost every $a$,
\begin{enumerate}
\item $R_a$ is symmetric, middle-positive, separately $d'$- and
$d''$-closed, coefficient-Radon, and $s$-annihilated;
\item its coefficient measures are carried by $s^{-1}(a)$;
\item for every Borel $E\subset J$ and every original compact
smooth test $\eta$,
\[
               \langle R_E,\eta\rangle
                    =\int_E\langle R_a,\eta\rangle\,d\nu(a).
 \tag{DI3}
\]
\end{enumerate}
Local coefficient variations of the family are $\nu$-integrable.
The exceptional set can be chosen once for all original tests and
all actual moment refinements. For the chosen control measure $\nu$,
the family is unique up to one $\nu$-null set as original currents.
At every parameter atom of positive $\nu$-mass,
\[
                       R_a=\nu(\{a\})^{-1}R_{\{a\}}.
 \tag{DI4}
\]
The normalization measure is not claimed to be canonical; the
closed-band current-valued measure $E\mapsto R_E$ is canonical.
\end{theorem}

\noindent\textbf{Status: \texttt{NEW\_PROVED}} for original
all-test currents. This is not an auxiliary-field current theorem.

\begin{proof}
Choose the countably cofinal stages of
Lemma~\ref{lem:r12-countably-cofinal-original-tests}. On each stage,
use a finite constant exterior coefficient frame with a fixed
compact cutoff. The coefficient-Radon property gives finitely many
finite signed Radon measures representing all its test pairings.
List the resulting countable collection as $(\mu_j)_{j\ge1}$,
including finitely many positive domination controls on every stage.
The measure
\[
            \mu=\sum_{j\ge1}
                    \frac{2^{-j}}{1+\|\mu_j\|_{\mathrm{TV}}}|\mu_j|
 \tag{DI5}
\]
is finite positive Radon and dominates every listed coefficient
measure by a fixed constant depending on its index. Thus every
fixed compact coefficient frame is bounded by $C\mu$, not merely
absolutely continuous without quantitative control.
If $R=0$ the theorem is trivial; otherwise discard any
zero summands and put $\nu=s_*\mu$.

Write $\mu_j=g_j\mu$ using Borel versions of their Radon--Nikodym
densities, and disintegrate $\mu$ by the preceding lemma. Outside
one $\nu$-null set all $\int|g_j|\,d\mu_a$ are finite. Put
$\mu_{j,a}=g_j\mu_a$. On a given stage these finitely many signed
Radon measures define an order-zero functional on its smooth
coefficient arrays.

We check that these functionals give an original current, rather
than independent unrelated functions on the chosen stages. For
each of the countably many dense zero-form relations from the
cofinality lemma, its original coefficient measure is zero, hence
its conditional pairing is zero almost everywhere. Impose all
these equalities outside one null set. Each conditional stage
functional is continuous for the smooth coefficient topology,
by its finite coefficient-variation bound, so it kills the entire
closed relation kernel. The same reasoning on common stages
identifies any two representations. Consequently the functionals
extend to one coefficient-Radon original current $R_a$ on all
compact original smooth tests. The measure construction proves
(DI3), first on the countable stages and then on every test.
Its finite stage variations are integrable because
\[
                 \int_J\int|g_j|\,d\mu_a\,d\nu
                                           =\int|g_j|\,d\mu<\infty.
 \tag{DI6}
\]

For every positive test in the chosen countable dense family,
positivity of every $R_E$ in (DI3) forces
$\langle R_a,\eta\rangle\ge0$ almost everywhere. Impose these
countably many inequalities simultaneously; continuity on each
stage gives positivity on its full middle-positive cone. Thus
$R_a$ is middle-positive, not merely positive on decomposable or
strongly positive tests. The same countable argument applied to
the symmetry identities gives symmetry.

Every $R_E$ is separately closed by the new closed-band theorem.
For each test $\alpha$ in a countable dense source stage,
(DI3) applied to $d'\alpha$ gives
\[
                      \int_E\langle R_a,d'\alpha\rangle\,d\nu=0
                           \qquad\text{for all Borel }E.
\]
Its integrand is therefore zero almost everywhere. Impose the
countably many prime and double-prime equations and extend by
continuity of the differentials on each source stage. This proves
separate closedness on all original tests. The identical argument
for $d's_i\wedge R_E$ and $d''s_i\wedge R_E$ proves (PB2) for
$R_a$. The support assertion follows from that of $\mu_a$.
Set $R_a=0$ on the one exceptional null set if a family indexed by
every parameter is desired; no potential function is changed.

Uniqueness follows from (DI3) on a countable determining family:
two families have equal pairings there almost everywhere, and
continuity then identifies their original currents. Applying (DI3)
to $E=\{a\}$ proves (DI4). A null exceptional set cannot contain
an atom of positive $\nu$-mass. All refinements were included in
the exact cofinality lemma, so the conclusion is not a statement
only on the initial countable moment names.
\end{proof}

\paragraph{The scope of the conditional object.}
Each $R_a$ is an original $K$-current on $U$, tested by the entire
ambient original sheaf, including moments with genuine mixed
dependence on any parameter and fiber coordinates. The theorem
does not identify it with a current on $Y$, on a norm-level
polyhedron, or on a completed-residue-field analytic fiber.
It asserts no inverse to a geometric pushforward. Its use in the
next section is legitimate precisely because the frozen BT theorem
already permits every original positive separately closed current
on $U$, including these conditional ones.
